\section{总结与延伸}
Attention alternatives 和 MoE 是现代 LLM 架构里最重要的两类非 dense 思路。前者改变 token 如何访问上下文，后者改变 token 如何访问参数。二者都能带来 scale 或 inference efficiency 的收益，但都把问题转移到 routing、state design、load balancing、communication 和 stability 上。
\subsection{拓展阅读}
\begin{longtable}{p{0.27\textwidth}p{0.31\textwidth}p{0.32\textwidth}}
\toprule
阅读对象 & 带着什么问题读 & 对应本讲证据 \\
\midrule
Linear attention、Mamba-2、Gated DeltaNet & 状态更新如何避免 $n^2$ matrix？哪些长程信息会丢失？ & Slides 4--10 的公式、recurrent form 和 hybrid performance。 \\
DeepSeek Sparse Attention & 重要位置如何被选择，稀疏索引开销是否抵消计算节省？ & Slides 11--13 的 sparse selection 与质量/速度结果。 \\
Switch Transformer、GShard & 经典 top-$k$ routing、capacity factor 与 load-balancing loss 如何形成？ & Slides 15--31 的 MoE 动机与 routing 基线。 \\
DeepSeek MoE、OlMoE、Mixtral、Qwen MoE & Fine-grained/shared experts、auxiliary-free balancing 与 upcycling 在什么条件下有效？ & Slides 32--53 的消融、系统和适配证据。 \\
DeepSeek v2/v3 & MoE、MLA 与 MTP 分别节省哪种资源，组合后新增什么耦合？ & Slides 54--59 的完整 case study。 \\
JAX Scaling Book & 如何把 attention/MoE 公式转换为 FLOPs、GPU 显存读写、collective 与 latency 账本？ & 本讲资源验收表和最小复现实验。 \\
\bottomrule
\end{longtable}

\begin{importantbox}{课后练习：为 selective computation 写一份反事实报告}
选择一个 attention alternative 或 MoE 配置，先写出它声称节省的资源，再构造一个可能让它更慢或更差的反事实环境。例如：短序列使 sparse indexing 得不偿失；小 batch 让 expert GEMM 太碎；慢互联让 all-to-all 主导；长程 retrieval 需要更高 full-attention 频率。最后列出能区分“机制无效”和“实现不佳”的指标。这个练习的目标不是支持某种架构，而是训练在部署前主动寻找失效边界。
\end{importantbox}
\end{document}
